\section{A prioritized research agenda}
\label{sec:research-questions}

The next objective is to turn the proved query $x\mapsto\D(x)$ into a
complete recognition procedure with an explicit global complexity bound.
The questions below separate near-term implementation work from structural
estimates that remain to be proved. Throughout, $n$ denotes the crossing
number of the supplied diagram, not the minimum crossing number of its knot.
Each proposed advance should identify its input, its checked output, and the
quantity that decreases or pays for its cost.

\subsection{First priority: certify the diagram-to-exterior construction}

Can the fast diagram pipeline produce a finite triangulation together with
efficiently checkable evidence that it represents the exterior of exactly
the input knot? A useful certificate should describe the local crossing and
arc blocks, their face identifications, and a peripheral marking. Subsequent
simplification should carry a sequence of verified local transformations or
another explicit homeomorphism certificate. The construction must have a
proved polynomial bound on tetrahedra and certificate length in the diagram
encoding.

This bridge makes a positive disc query relevant to the original recognition
problem. The immediate experiment is to attach provenance to existing
diagram-derived fixtures and compare the resulting finite triangulations
with an independent topology implementation. Such comparisons guide
development; the mathematical obligation is a proof for every accepted
construction. Mutation tests should change an arc attachment, a peripheral
label, or a simplification step and require the checker to detect the broken
correspondence.

\subsection{First priority: a complete constrained disc producer}

What search procedure supplies the normal vectors, and how many must it
inspect? The delivered count eliminates a candidate with $\D(x)=0$; it gives
no exhaustion certificate for the remaining admissible vectors. Develop a
producer with a clearly specified finite search space, exact feasibility
bounds, and independently justified pruning rules. Record search nodes,
candidate bit lengths, and time spent in feasibility, lifting, and component
queries separately.

An essential disc can already be found on a Q-vertex ray under the hypotheses
of Suh's theorem; proving that existence again is not the research question
\cite{Suh2010}. The difficult question is whether a complete search can avoid
exponentially many relevant branches on every input. The established
branch-and-bound framework provides a practical comparison
\cite{BurtonOzlen2014}. A proposed pruning rule needs a proof that some
essential-disc witness survives whenever one exists, including when an
intermediate candidate is disconnected or has nonorientable components.

\subsection{First priority: canonicalize before expensive search operations}

Can candidate generation operate directly on primitive quadrilateral data,
constructing the canonical standard lift only when required? Canonical
lifting and vertex-link ambiguity are established Q-theory
\cite{Tollefson1998,Burton2009}. The opportunity is to use them earlier in the
implemented search, so that large common multiplicities and independent
vertex-link offsets never enter its costly orbit queries.

For every proposed Q-coordinate input, construct or certify its integral
lift and check the full admissibility conditions. Evaluate whether branching
bounds can use the coordinate estimate of
Theorem~\ref{thm:primitive-q-core}. The first experiment should compare
search-node and arithmetic counts before and after canonicalization on the
same candidate stream. The identity $\D(x)=g\D(P)$ supports reuse of an
essential-disc count. A request for the original component histogram still
requires its original surface data: vertex-link deletion changes that
histogram, and scaling a one-sided component changes its sheet decomposition.

\subsection{Near-term refinement: certify and benchmark the two-weight mode}

The proposition in \cref{sec:two-weight-refinement} already proves that one
ambient-certified boundary probe suffices for all supplied surfaces in a
fixed supported manifold. Implement its finite binary preprocessing and its
linear-size cocycle witness, then compare two and three weights on identical
candidate streams. The open performance question is the break-even point
between additional ambient preprocessing and the saved transport coordinate.
Cache only source-verified preprocessing data, and include cold and amortized
costs separately. The existing three-weight implementation remains the
reference until the new mode passes the frozen external corpus and the
source-mutation checks. This is an implementation and cost question, not a
conjecture about whether the two-weight criterion is true.

\subsection{Second priority: batch related candidates and reuse checked traces}

How much computation can be shared across candidates in one triangulation?
The canonical core gives an exact starting point: candidates
$x_j=\sum_v\ell_{jv}L_v+g_jP$ share the answer $\D(P)$, with their counts
recovered by multiplication. Cache the validated triangulation, canonical
cores, and their certificates using explicit source identifiers. Charge the
cost of reading each candidate and checking its reduction even on cache hits.

A broader question is whether an interval-orbit trace can be reused over a
region of coordinate space. Equal quadrilateral support alone does not fix
endpoint order or scheduler decisions. A reusable trace would need certified
inequalities and congruences defining where every step remains valid, including
periodic cases. The completed replay artifacts provide a test harness for
this experiment \cite{AHT2006}. Measure reuse per distinct core and per
verified trace region, and retain the total candidate count in the global
cost accounting.

\subsection{Second priority: connect selected components to compressed cutting}

Can a positive three-weight query return enough additional information to
cut along one certified essential disc efficiently? Recover a selected
component's normal vector, using coordinate weights or a targeted orbit
query, and supply the gluing and boundary data needed by the next stage.
This tests a concrete transition from a component predicate to a geometric
operation. A histogram with positive disc multiplicity establishes existence,
but does not by itself specify the cut triangulation.

Compressed cutting must preserve the hypotheses of the routine being used.
Lackenby's hierarchy framework imposes explicit irreducibility and essential
boundary-pattern conditions \cite{Lackenby2026}. Determine how each condition
is established initially and propagated through subsequent cuts, and include
the relevant boundary-pattern evidence in each transition certificate.
Efficient algorithms for curves and normalization offer useful subroutines
\cite{LackenbyCurves2026}; the surrounding three-dimensional invariants still
require proof. Begin with cuts whose outputs have small independent finite
oracles before enlarging the binary weights.

\subsection{Structural priority: prove an amortized hierarchy budget}

Which structural invariant can limit the aggregate search in an accelerated
hierarchy? For a proposed encoding with counters ranging from zero to $g_i$,
a useful \emph{target}, not a theorem established here, is
\[
                 \sum_i\log_2(g_i+1)=O(\log^2 n).
\]
These $g_i$ are hierarchy counter bounds, unrelated to the quadrilateral gcd
of a single input. If all other state choices are charged within the same
budget, the product of the counter ranges is
$2^{O(\log^2 n)}=n^{O(\log n)}$.

The proposed bound must control all reachable states relevant to the
algorithm, including branching and counter resets. A bound along one chosen
successful path is insufficient for an exhaustive search. Seek an invariant
that pays for each split, new boundary pattern, and increase in a counter's
range. The explicit hierarchy machinery in \cite{Lackenby2026} supplies a
precise setting in which to formulate this obligation. Failure to obtain the
target should produce an explicit obstruction family or a weaker proved
budget, rather than an empirical asymptotic claim.

\subsection{Structural priority: anisotropic bounds and a graph of states}

Can individual counter bounds replace a uniform worst-case cap? If an
encoding has $m$ counters, retain $\prod_{i=1}^{m}(g_i+1)$ instead of replacing
it immediately by $(1+\max_i g_i)^m$. This mixed-radix accounting can expose
instances with many small choices and only a few large ones. Identify a
topological reason for that anisotropy and test whether it persists after
cutting and simplification.

Can histories reaching the same certified state be merged into a directed
acyclic graph? State equivalence must preserve the manifold, boundary
pattern, peripheral information, and remaining search obligations. Prove
acyclicity using a well-founded rank, or account explicitly for cycles.
Canonical labels and hashes may locate possible matches, but each merge
needs a verified equivalence. Initial measurements should count distinct
states, transitions, and equivalence-checking cost. A small state graph is
useful only when constructing and recognizing its states is also affordable.

\subsection{Quantitative priority: an explicit deterministic bit bound}

What exponent does the composed algorithm actually satisfy? Derive separate
bounds for triangulation construction, candidate production, integer gcds,
interval operations, weight-run growth, component recovery, cutting, and
certificate checking. Track maximum operand bit lengths and the cost
$M(b)$ of multiplying $b$-bit integers. The weighted AHT theorem gives a
polynomial local foundation \cite{AHT2006}; composition still requires bounds
on the number and size of its invocations.

Publish a deterministic bound with fixed constants, distinguishing
$2^{O((\log n)^2)}$ from $2^{O((\log n)^3)}$ and stating which exponent is
actually proved. Include output and certificate lengths, memory, and input
reading. The completed operation counters and paired kernel measurements
can identify costly terms and check predictions. Their role is to guide
which inequalities need improvement; the final worst-case bound must follow
from the algorithm and its invariants.

\subsection{Continuous priority: stronger adversarial evidence}

Extend the delivered 1,275-vector, 45-triangulation Regina corpus and its
isolated counterexamples. Preserve both essential-disc cancellation and the
orientable sphere-plus-punctured-torus false positive as mandatory cases.
Add large vertex valence, additional boundary vertices, long layered
triangulations, coprime large quadrilateral coordinates, and inputs designed
to increase interval subdivisions. Exercise one-sided scaling and certificate
mutations at every new interface.

Keep a capped corpus with fully expanded independent components and a
separate collection of large binary instances whose answers follow from
proved constructions. For each family, report the parameter varied, the
expected component semantics, the source triangulation, and reproducible
seeds. When recognition stages are connected, add difficult diagram families
and measure every stage on the same instances. This prevents a local speedup
from hiding a newly dominant producer or source-validation cost, while
maintaining the exact distinction between a component of a supplied surface
and a conclusion about the original knot.
